% ECONOMICS_PROBLEM: {"id": "L41-109", "title": "Algorithmic pricing and collusion", "jel_code": "L41", "source_id": "OP-0074"}
% !TeX program = xelatex
\documentclass[11pt,a4paper]{article}
\usepackage[margin=23mm]{geometry}
\usepackage{fontspec}
\setmainfont{Latin Modern Roman}
\usepackage{amsmath,amssymb,mathtools}
\usepackage{xcolor,enumitem,booktabs,longtable,array,xurl,hyperref}
\definecolor{muted}{HTML}{53636C}
\hypersetup{colorlinks=true,urlcolor=blue,linkcolor=blue}
\setlength{\parindent}{0pt}
\setlength{\parskip}{5pt}
\newcommand{\R}{\mathbb R}
\newcommand{\E}{\mathbb E}
\newcommand{\N}{\mathbb N}
\newcommand{\ind}{\mathbf 1}
\newcommand{\Var}{\operatorname{Var}}
\newcommand{\Cov}{\operatorname{Cov}}
\newcommand{\supp}{\operatorname{supp}}
\DeclareMathOperator*{\argmin}{arg\,min}
\DeclareMathOperator*{\argmax}{arg\,max}
\newcommand{\entrymeta}[3]{{\sffamily\small\color{muted}\textbf{\texttt{#1}}\quad|\quad #2\par #3\par}}
\newcommand{\Problemref}[1]{\texttt{#1}}
\newcommand{\JELref}[2]{\texttt{#2} (JEL \texttt{#1})}
\begin{document}
\section*{Algorithmic pricing and collusion}
\textbf{Problem ID:} \texttt{L41-109}\quad \textbf{JEL:} \texttt{L41}\par
% BEGIN ECONOMICS STATEMENT
\label{problem:OP-0074}
\label{atlas:prob:I4}

\noindent\textbf{Original identifier:} I4 (atlas item 74). \textbf{Classification:} Algorithmic competition robustness research program. \textbf{Original tractability label:} Medium (editorial, not a complexity result).

\subsection*{Definitions and assumptions}

Consider a specified repeated differentiated-product price game with firms $i=1,\ldots,n$, demand $q_i(p,s)$, marginal costs $c_i$, and common state $s$. Each independently installed algorithm maps its own information history and internal learning state into a price distribution. Define a benchmark $p^N(s)$ as a selected one-shot Nash equilibrium of the same demand and cost game; competition need not mean zero markups. Define average excess profit or price relative to this benchmark, weighted by prespecified quantities. Specify information sharing, vendor overlap, exploration, replacement, entry, and discounting. Independence of installation does not imply statistical independence of observed prices.

\subsection*{Formal research question}

Characterize algorithm and market classes for which supracompetitive behavior is robust. For example, with $G_t$ the declared profit gap above the benchmark, determine conditions yielding
\[
\Pr\!\left(\liminf_{T\to\infty}\frac1T\sum_{t=1}^TG_t\geq\varepsilon\right)\geq1-\alpha
\]
for fixed $\varepsilon>0$ and $\alpha\in(0,1)$, or establish finite-horizon analogues when learning does not converge. Evaluate robustness to heterogeneous vendors and profitable unilateral deviations. Separately identify causal effects of adoption on market outcomes using randomized adoption or credible exogenous variation, with explicit assumptions about simultaneous adoption and spillovers.

\subsection*{Interpretation and scope}

The bare existence question has already received affirmative constructive evidence in particular simulations. The substantive extension concerns when that behavior survives realistic costs, demand shocks, algorithm turnover, and entry. Parallel prices can reflect common costs or competitive response and therefore do not identify tacit cooperation. Nor does an above-Nash average price prove legal collusion: the economic benchmark and legal agreement standard are distinct. A learning algorithm may fail to converge, so a limiting-price claim requires a convergence proof or a justified stochastic stationary distribution. Finite-run simulation estimates should report sensitivity to random seeds, initialization, and exploration schedules and cannot prove the displayed probability bound for all environments. Deviation tests need to allow firms to modify algorithms, not merely change one price while retaining a fixed learning rule. Empirical adoption studies additionally confront endogenous timing and local market exposure. A model-based robustness result and a causal market estimate are complementary deliverables with different evidential requirements.

\subsection*{Source audit and status}

[S1] provides Q-learning simulation evidence; [S2], published in 2024, provides German gasoline-market adoption evidence; [S3] analyzes legal-economic issues involving autonomous agents. All citations are verified. Consequently this item is corrected from an unanswered existence claim to a robustness and identification frontier. The displayed asymptotic criterion is an editorial target and is not asserted as a theorem or universal empirical fact established by these papers.

\subsection*{References}

\begin{enumerate}[label={[S\arabic*]},leftmargin=*]

\item Emilio Calvano, Giacomo Calzolari, Vincenzo Denicolò, and Sergio Pastorello (2020). \emph{Artificial Intelligence, Algorithmic Pricing, and Collusion}. American Economic Review 110(10), 3267--3297. \url{https://www.aeaweb.org/articles?id=10.1257/aer.20190623}\par \textit{Source role:} Constructive simulation evidence for algorithmic supracompetitive prices. \textit{Locator:} Abstract and article metadata.

\item Stephanie Assad, Robert Clark, Daniel Ershov, and Lei Xu (2024). \emph{Algorithmic Pricing and Competition: Empirical Evidence from the German Retail Gasoline Market}. Journal of Political Economy 132(3), 723--771. \url{https://doi.org/10.1086/726906}\par \textit{Source role:} Market-specific evidence on adoption and competition. \textit{Locator:} Abstract and article metadata.

\item Joseph E. Harrington, Jr. (2018). \emph{Developing Competition Law for Collusion by Autonomous Artificial Agents}. Journal of Competition Law \& Economics 14(3), 331--363. \url{https://doi.org/10.1093/joclec/nhy016}\par \textit{Source role:} Legal-economic conceptual analysis. \textit{Locator:} Abstract and article metadata.

\end{enumerate}

\subsection*{Common conventions: Source mathematical conventions}

Unless an entry states otherwise, agent and firm sets are finite. State and action spaces are measurable, strategies and outcome maps are measurable, probability laws are specified, and expectations exist for the targets under discussion. Infinite-horizon discount factors lie in $(0,1)$ and discounted objectives are finite. Norms, tolerances and complexity input sizes must be specified for computational comparisons. Constants or primitives that appear locally are inputs, not empirical facts asserted by this catalogue.

\textbf{Identification.} A statistical model $m\in\mathcal M$ implies an observable law $P_m$ and target $\theta(m)$. At law $P$, its identified set is
\[
\Theta(P;\mathcal M)=\{\theta(m):m\in\mathcal M,\ P_m=P\}.
\]
Point identification holds when this set is a singleton, or equivalently when $P_{m_1}=P_{m_2}$ implies $\theta(m_1)=\theta(m_2)$ for all compatible models. Sharp bounds mean bounds attained, or approached when the boundary is not attained, by compatible models. Writing this set is a definition, not a solution: the substantive task is to establish informative restrictions and derive or estimate the set. An unrestricted-confounding model may give only trivial bounds.

\textbf{Inference.} Identification, estimability and valid uncertainty quantification are separate questions. Fix $\alpha\in(0,1)$. A confidence procedure $C_n$ over a declared law class $\mathcal P$ has asymptotic uniform coverage at level $1-\alpha$ when
\[
\liminf_{n\to\infty}\inf_{P\in\mathcal P}\Pr_P\{\theta(P)\in C_n\}\geq1-\alpha.
\]
For set-identified targets the coverage event must be defined separately, for example coverage of the full identified set. If an entry uses identification-indexed models instead of uniquely indexed laws, the same coverage convention is applied over those models. Pointwise limits do not establish uniform validity, and asymptotic validity does not guarantee finite-sample precision. Sample sizes, dependence structures and weak-identification sequences are part of each application.

\textbf{Counterfactuals.} $Y_i(a)$ is a potential outcome under a specified intervention, including its timing, implementation and versions. It is not $Y_i$ conditional on voluntarily choosing $a$. A full intervention vector permits interference; shorthand scalar treatment notation does not silently impose no interference. Assignment, selection, attrition, anticipation and measurement must be specified. A claim about an instrument's local effect is not automatically a population policy effect.

\textbf{Economic equilibrium.} A counterfactual model includes best responses, constraints, feasibility, market clearing, state transitions and expectations consistent with the stated information. When several equilibria exist, either a declared selector is an input or the target is set-valued. Comparisons across policy regimes must use the stated selection convention. Reduced-form relationships are not presumed invariant after an equilibrium-changing intervention.

\textbf{Optimization and welfare.} Optimization is over the stated feasible set. If attainment is not established, $\sup$ or $\inf$ and $\varepsilon$-optimal choices replace an asserted maximizer or minimizer. For example, with value $V=\sup_{a\in\mathcal A}W(a)<\infty$, an $\varepsilon$-optimal set is $\{a\in\mathcal A:W(a)\geq V-\varepsilon\}$ for $\varepsilon>0$. Benefits, costs, risk and health or environmental outcomes can be aggregated only using declared units and conversion weights. Interpersonal utility weights, the treatment of fiscal transfers and foreign profits, discounting, and the behavioral welfare criterion are explicit inputs. A comparative-static derivative additionally requires existence and appropriate differentiability of the selected counterfactual.

\textbf{Sources.} Local labels [S1], [S2], and so forth restart in every question. Locators identify the relevant abstract, model, section or result at the level actually checked; a bibliographic or abstract-level verification is not represented as inspection of an entire proof. Working papers are distinguished from later publications and changing versions. Source summaries are paraphrased. All URL checks and editorial formulations refer to the audit date stated above.
% END ECONOMICS STATEMENT
\end{document}
